\section{Résultats de l'étape 1}

\subsection{Benchmark compute-only}

\begin{table}[H]
    \centering
    \begin{tabular}{c c c c}
        \toprule
        \textbf{Threads} & \textbf{Temps moyen par pas (ms)} & \textbf{Speed-up} & \textbf{Efficacité} \\
        \midrule
        1 & 227.970 & 1.000 & 1.000 \\
        2 & 124.934 & 1.825 & 0.912 \\
        4 & 69.480 & 3.281 & 0.820 \\
        8 & 48.834 & 4.668 & 0.584 \\
        16 & 33.233 & 6.861 & 0.429 \\
        \bottomrule
    \end{tabular}
    \caption{Résultats du benchmark compute-only pour l'étape 1}
    \label{tab:stage1_results}
\end{table}

\subsection{Graphiques du benchmark compute-only}

\begin{figure}[H]
    \centering
    \begin{subfigure}[t]{0.48\linewidth}
        \centering
        \begin{tikzpicture}
\begin{axis}[
    width=\linewidth,
    height=6.0cm,
    xlabel={Nombre de threads},
    ylabel={Speed-up},
    xmin=1,
    xmax=16,
    ymin=0,
    xtick={1,2,4,8,16},
    grid=both,
    minor y tick num=1,
    axis line style={draw=black!60},
]
\addplot[
    color=enstaBleuFonce,
    mark=*,
    line width=1.1pt,
]
coordinates {
        (1, 1.000)
        (2, 1.825)
        (4, 3.281)
        (8, 4.668)
        (16, 6.860)
};
\end{axis}
\end{tikzpicture}

        \caption{Speed-up}
        \label{fig:stage1_speedup}
    \end{subfigure}
    \hfill
    \begin{subfigure}[t]{0.48\linewidth}
        \centering
        \begin{tikzpicture}
\begin{axis}[
    width=\linewidth,
    height=6.0cm,
    xlabel={Nombre de threads},
    ylabel={Efficacité},
    xmin=1,
    xmax=16,
    ymin=0,
    ymax=1.05,
    xtick={1,2,4,8,16},
    grid=both,
    minor y tick num=1,
    axis line style={draw=black!60},
]
\addplot[
    color=enstaBleuClair,
    mark=square*,
    line width=1.1pt,
]
coordinates {
        (1, 1.000)
        (2, 0.912)
        (4, 0.820)
        (8, 0.584)
        (16, 0.429)
};
\end{axis}
\end{tikzpicture}

        \caption{Efficacité}
        \label{fig:stage1_efficiency}
    \end{subfigure}
    \caption{Scalabilité observée sur le benchmark compute-only de l'étape 1}
    \label{fig:stage1_scaling}
\end{figure}

\subsection{Benchmark interactif}

Pour préparer la comparaison avec l'étape 2, on a aussi mesuré l'étape 1 en mode interactif, avec fenêtre ouverte, sur \texttt{data/galaxy\_1000}. Le temps total correspond ici à \texttt{Render time + Update time}.

\begin{table}[H]
    \centering
    \resizebox{\linewidth}{!}{%
    \begin{tabular}{c c c c c c}
        \toprule
        \textbf{Threads} & \textbf{Render (ms)} & \textbf{Update (ms)} & \textbf{Total (ms)} & \textbf{Speed-up total} & \textbf{Efficacité totale} \\
        \midrule
        1 & 4.500 & 16.800 & 21.300 & 1.000 & 1.000 \\
        2 & 5.025 & 9.150 & 14.175 & 1.503 & 0.751 \\
        4 & 4.750 & 5.800 & 10.550 & 2.019 & 0.505 \\
        8 & 4.475 & 4.300 & 8.775 & 2.427 & 0.303 \\
        16 & 4.688 & 3.812 & 8.500 & 2.506 & 0.157 \\
        \bottomrule
    \end{tabular}%
    }
    \caption{Résultats du benchmark interactif pour l'étape 1}
    \label{tab:stage1_interactive}
\end{table}

\begin{figure}[H]
    \centering
    \begin{subfigure}[t]{0.48\linewidth}
        \centering
        \begin{tikzpicture}
\begin{axis}[
    width=\linewidth,
    height=6.0cm,
    xlabel={Nombre de threads},
    ylabel={Speed-up total},
    xmin=1,
    xmax=16,
    ymin=0,
    xtick={1,2,4,8,16},
    grid=both,
    minor y tick num=1,
    axis line style={draw=black!60},
]
\addplot[
    color=enstaBleuFonce,
    mark=*,
    line width=1.1pt,
]
coordinates {
        (1, 1.000)
        (2, 1.503)
        (4, 2.019)
        (8, 2.427)
        (16, 2.506)
};
\end{axis}
\end{tikzpicture}

        \caption{Speed-up total}
        \label{fig:stage1_interactive_speedup}
    \end{subfigure}
    \hfill
    \begin{subfigure}[t]{0.48\linewidth}
        \centering
        \begin{tikzpicture}
\begin{axis}[
    width=\linewidth,
    height=6.0cm,
    xlabel={Nombre de threads},
    ylabel={Efficacité totale},
    xmin=1,
    xmax=16,
    ymin=0,
    ymax=1.05,
    xtick={1,2,4,8,16},
    grid=both,
    minor y tick num=1,
    axis line style={draw=black!60},
]
\addplot[
    color=enstaBleuClair,
    mark=square*,
    line width=1.1pt,
]
coordinates {
        (1, 1.000)
        (2, 0.751)
        (4, 0.505)
        (8, 0.303)
        (16, 0.157)
};
\end{axis}
\end{tikzpicture}

        \caption{Efficacité totale}
        \label{fig:stage1_interactive_efficiency}
    \end{subfigure}
    \caption{Graphiques du benchmark interactif pour l'étape 1}
    \label{fig:stage1_interactive_scaling}
\end{figure}

On observe une amélioration nette du temps total, mais plus limitée qu'en compute-only. La partie \texttt{Render time} reste presque constante, ce qui borne mécaniquement le gain global.

\begin{figure}[H]
    \centering
    \begin{tikzpicture}
\begin{axis}[
    ybar,
    bar width=22pt,
    width=0.82\linewidth,
    height=6.0cm,
    ymin=0,
    ylabel={Temps (ms)},
    symbolic x coords={Render, Update},
    xtick=data,
    nodes near coords,
    every node near coord/.append style={font=\footnotesize},
    enlarge x limits=0.35,
    grid=both,
    minor y tick num=1,
    axis line style={draw=black!60},
]
\addplot[fill=enstaBleuClair, draw=enstaBleuFonce] coordinates {
    (Render, 4.750)
    (Update, 5.800)
};
\end{axis}
\end{tikzpicture}

    \caption{Décomposition des temps interactifs de l'étape 1 à 4 threads}
    \label{fig:stage1_interactive}
\end{figure}
